% =====================================================================
% §2.3 — La posizione di lettura (esempio: ex2_pointer)
%
% Pattern: domanda musicale -> diff YAML (un solo blocco) ->
%          fallimento della waveform -> PRIMA figura-partitura con
%          legenda minima -> meccanismo -> artefatto del freeze ->
%          ponte a §2.4
%
% TODO prima della submission:
%   [ ] linerange del listato: allineare alle righe reali di
%       examples/ex2_pointer/ex2_pointer.yml (mostrare solo il blocco pointer)
%   [ ] verificare nome figura: ex2_pointer_score.pdf
%   [ ] verificare il punto di congelamento (~1,3 s) dopo la taratura
%       definitiva dell'envelope, e che cada su vocale
%   [ ] \cite{Lippe1993cim}: allineare alla chiave reale in refs.bib
% =====================================================================

% =====================================================================
% NOTA WIP (ripristino su richiesta): di questa sottosezione esistono DUE
% versioni, entrambe con \label{sec:pointer}, tenute insieme di proposito
% per la fusione. NON e' il doppione gia' risolto: qui sono volutamente
% entrambe. In questo stato il PDF NON compila pulito (label e lst:pointer
% doppi; la v1 usa la figura ex2_pointer_log, assente). Fondere le due.
% =====================================================================

% --- v1 (logging grafico; figura ex2_pointer_log; con la formula integrale) ---
% ----------------------------------------------------------------
\subsection{La posizione di lettura}\label{sec:pointer}

In \S\ref{sec:c-e} il puntatore alla lettura 
è pilotato da un fasore il cui periodo coincide
con la durata del file audio di partenza. Il rapporto
tra i due è la \texttt{speed\_ratio} ($\tfrac{T_{\phi}}{D}$), 
che nell'esempio \texttt{ex2\_pointer.yml}  (Listato~\ref{lst:pointer}) viene espressa
come un \texttt{Envelope}\footnote{Un \texttt{Envelope} è una spezzata --- una funzione lineare a tratti, data
per breakpoint $[t, v]$ e interpolata fra l'uno e l'altro; con
\texttt{time\_mode: normalized} l'ascissa $t$ corre da $0$ a $1$ come frazione
della durata dello stream. Ogni parametro accetta indifferentemente uno scalare
o un envelope, e d'ora in poi la sintassi si dà per nota}.

\lstinputlisting[
      linerange={34-}, % TODO: allineare al file reale dopo l'aggiornamento dell'envelope
      language=yaml,
      basicstyle=\ttfamily\fontsize{7.5}{9}\selectfont,
      label=lst:pointer,
      caption={L'unica chiave aggiunta alla trasformazione identica
      (\texttt{ex2\_pointer.yml}).}
    ]{examples/ex2_pointer/ex2_pointer.yml}

 Quello dell'esempio
non è monotòno: \texttt{speed\_ratio} parte da $1$ --- dove lo stream è ancora
la riproduzione fedele di \S\ref{sec:c-e} --- sale a $2{,}3$, ridiscende oltre
lo zero fino a $-0{,}3$ e da ultimo si arresta a $0$, dove rimane sino alla
fine.

La trasformazione sfugge alla forma d'onda --- l'ampiezza resta piena e
continua, ché a mutare non è quanta energia scorre ma quale punto del segnale
la genera. Occorre una rappresentazione della lettura, ed è il logging grafico
che il sistema emette a ogni rendering (Fig.~\ref{fig:pointer-log}): in ascissa
il tempo del brano, in ordinata la posizione di lettura nel buffer sorgente,
ogni segmento un grano. La traiettoria dichiarata vi si legge per intero --- la
salita convessa mentre il percorso accelera in avanti, il colmo dove la
velocità attraversa lo zero, la discesa mentre la lettura si inverte, e
l'orizzontale del congelamento in coda al brano. Al logging grafico, qui
ridotto alla sua legenda minima, è dedicata la \S\ref{sec:partitura}.

\begin{figure}[h]
    \includegraphics[width=\linewidth]{examples/ex2_pointer/ex2_pointer_score}\\
    \captionof{figure}{Logging grafico del rendering di \texttt{ex2\_pointer}:
    ascisse = tempo del brano, ordinate = posizione di lettura nel buffer
    sorgente ($0$ = inizio, $1$ = fine del segnale); ogni segmento un grano. Il
    puntatore accelera in avanti, raggiunge un colmo, legge a ritroso e infine
    si congela.}
    \label{fig:pointer-log}
\end{figure}

Il principio è l'integrazione. La posizione di lettura è l'accumulo della
velocità sul tempo del brano,
\[
  \mathrm{pos}(t) \;=\; \texttt{start} \;+\; \int_0^t \texttt{speed\_ratio}(u)\, du,
\]
sicché \texttt{speed\_ratio} ne è la velocità istantanea e la posizione la fase
accumulata --- un fasore che percorre il buffer. Il suo è un rapporto: la
velocità con cui il puntatore avanza, riferita a quella della lettura nativa.
Vale $1$ quando le due coincidono, $0$ quando il puntatore è fermo, sotto $1$
quando avanza più lento e sopra $1$ più rapido; col segno negativo il percorso
si inverte. Di qui i regimi dell'esempio --- $2{,}3$ una corsa in avanti,
$-0{,}3$ un'inversione lenta, $0$ il congelamento.

Un punto va isolato, ed è quello da cui dipende l'intera espressività
dell'asse. Ogni grano legge il buffer alla frequenza propria del segnale, e
l'altezza che porta è quella che il materiale ha nel punto da cui è tratto;
\texttt{speed\_ratio} dispone soltanto di quel punto e del passo con cui
avanza. Percorso e altezza sono perciò grandezze separate --- qui muoviamo il
percorso, e l'altezza, cui spetta un parametro proprio, resta quella del
segnale. È questa indipendenza a fare della posizione di lettura il parametro
espressivo dominante della granulazione~\cite{Lippe1993cim}: si può sostare,
accelerare, invertire, congelare lungo il materiale, e l'altezza resta
governata a parte.

La specifica è intanto ancora del tutto deterministica: nessun valore è estratto
a caso, e due rendering coincidono al campione. Alla fine del brano il puntatore
è ormai fermo, ma i grani che emette nascono ancora sulla griglia metronomica
ereditata dai default: è quella regolarità lo scostamento seguente.
